\begin{table}[H]
    \centering
    \caption{Результаты измерений для определения основной погрешности}
    \label{tab:error_measurement}
    \begin{tblr}{
        width = \linewidth,
        colspec = {X[c,m] X[c,m] X[c,m]},
        hlines, vlines,
    }
        {Поверяемая отметка \\ на ЭВ, $U_i$, В} & {Показание ЦВ при \\ увеличении, $U_{\text{о.ув},i}$, В} & {Показание ЦВ при \\ уменьшении, $U_{\text{о.ум},i}$, В} \\
        %
        % Заполнить 6-10 строк
        %
        & & \\
        & & \\
        & & \\
        & & \\
        & & \\
        & & \\
        & & \\
        & & \\
        & & \\
        & & \\
    \end{tblr}
    \textit{Нормирующее значение (верхний предел шкалы): $U_N = \_\_\_\_\_\_~\text{В}$}
\end{table}

\begin{longtblr}[
    caption = {Результаты измерений для области верхних частот},
    label = {tab:afc_high_freq},
]{
    width = \linewidth,
    colspec = {X[c,m] X[c,m]},
    hlines, vlines,
    rowhead = 1,
}
    {Частота, $f$, кГц} & {Показание ЭВ, $U(f)$, В} \\
    1 (опорная) & \\
    100 & \\
    200 & \\
    300 & \\
    400 & \\
    500 & \\
    600 & \\
    700 & \\
    800 & \\
    900 & \\
    1000 & \\
    1100 & \\
    1200 & \\
    1300 & \\
    1400 & \\
    1500 & \\
\end{longtblr}

\begin{longtblr}[
    caption = {Результаты измерений для области нижних частот},
    label = {tab:afc_low_freq},
]{
    width = \linewidth,
    colspec = {X[c,m] X[c,m]},
    hlines, vlines,
    rowhead = 1,
}
    {Частота, $f$, Гц} & {Показание ЭВ, $U(f)$, В} \\
    1000 (опорная) & \\
    800 & \\
    600 & \\
    400 & \\
    200 & \\
    100 & \\
    50 & \\
    40 & \\
    30 & \\
    20 & \\
    10 & \\
    8 & \\
    6 & \\
    4 & \\
    2 & \\
    0 & \\
\end{longtblr}

\begin{table}[H]
    \centering
    \caption{Показания вольтметра для сигналов различной формы}
    \label{tab:waveform_influence}
    \begin{tblr}{
        colspec = {Q[l,m] X[c,m]},
        hlines, vlines,
    }
        {Форма сигнала} & {Показание ЭВ, $U_\text{п}$, В} \\
        Синусоидальная & \\
        Прямоугольная & \\
        Треугольная & \\
    \end{tblr}
\end{table}